\documentclass[11pt,twoside]{article}
\usepackage{graphicx} % Required for inserting images
%% The amssymb package provides various useful mathematical symbols
\usepackage{amssymb}
%% The amsmath package provides various useful equation environments.
\usepackage{amsmath}
\usepackage{physics}
\usepackage[separate-uncertainty=true,group-minimum-digits=3, group-separator = {,}]{siunitx}
\usepackage[english]{babel}
\usepackage[utf8]{inputenc}
\usepackage[T1]{fontenc}
\usepackage{booktabs}
\usepackage{newtxtext,newtxmath}
\usepackage{csquotes} % Recommended for biblatex with babel/polyglossia
\usepackage{xcolor}
\usepackage{soul}
\usepackage{wasysym}
\usepackage[margin=0.75in]{geometry}
\usepackage{hyperref}
\hypersetup{
    colorlinks=true,
    linkcolor=blue,
    filecolor=magenta,      
    urlcolor=cyan,
    pdftitle={Mass loss model for pebble aggregate samples in DiMES},
    pdfpagemode=FullScreen,
    }
\usepackage{cancel}
\usepackage{listings}

\lstset{
    language=Python,
    frame=tBLR,
    basicstyle=\ttfamily\small,
    keywordstyle=\color{blue},
    commentstyle=\color{green!60!black},
    stringstyle=\color{red},
    showstringspaces=false,
    numberstyle=\tiny,
    numbers=left
}

\renewcommand{\lstlistingname}{Code}
\providecommand{\lstlistingautorefname}{Code}


\title{Mass loss model for pebble aggregate samples in DiMES}
\author{Erick Martinez Loran}
\date{August 2026}

\begin{document}

\maketitle

\section{The mass loss is proportional to the incident heat load}
\begin{align}
\frac{dm}{dt} = c q_{\mathrm{incident}} =  c q_{\parallel} \label{eq:mass-loss-q} A_w,
\end{align}
where $q_{\parallel}$ is the parallel heat load at the divertor and $A_w$ is the wetted area: $A_w(t) = Dh(t)/2$, with $h(t)$ the height of the rod exposed to the plasma at time $t$.

Consider the incident heat loss as
\begin{align}
q_{\parallel} &= \gamma n_e k_{\mathrm{B}} T_e c_s, \label{eq:plasma_heat_load}
\end{align}
where $\gamma$ is the transmission coefficient in the sheath, $k_{\mathrm{B}}$ is Boltzmann's constant, $c_s$ is the ion sound speed, $T_e$, and $n_e$ are the electron temperature and electron density at the sheath edge, respectively.

The ion saturation current density $J_{\mathrm{sat}}$ is the maximum current that can be drawn to a material surface (such as a Langmuir probe or vacuum chamber wall) and is defined as follows:
\begin{align}
J_{\mathrm{sat}} \equiv e n_{\mathrm{se}} c_s \label{eq:ion-saturation}
\end{align}

Substitute $n_{\mathrm{se}} c_s = J_{\mathrm{sat}} /e$ into the sheath transmission equation Eq.~\eqref{eq:plasma_heat_load}:
\begin{align}
q_{\parallel} &= \frac{\gamma J_{\mathrm{sat}} k_{\mathrm{B}} T_e}{e}. \label{eq:q_para_Jsat}
\end{align}

Multiply Eq.~\eqref{eq:q_para_Jsat} by the wetted area $A_w$ and plug it into Eq.~\eqref{eq:mass-loss-q}:
\begin{align}
\frac{dm}{dt} = c A_w \frac{\gamma J_{\mathrm{sat}} k_{\mathrm{B}} T_e}{e}
\end{align}

We propose that $A_w(t) J_{\mathrm{sat}} (t)$ is the current measured at DiMES:
\begin{align*}
\frac{dm}{dt} = c \gamma k_{\mathrm{B}} I(t)  T_e(t).
\end{align*}
The prefactor $c \gamma k_{\mathrm{B}}$ can be lumped into a single constant $c$
\begin{align}
\boxed{
\frac{dm}{dt} = c I(t) T_e(t) \label{eq:mass-loss-I}
}
\end{align}


To calibrate the model, we measure the total mass change $m_{\mathrm{L}} \equiv \Delta m = m_f - m_i$ before and after exposure to the plasma heat load.
Then,
\begin{align}
m_{\mathrm{L}} = c \int_0^t I(t)  T_e(t) dt = cQ. \label{eq:mass-loss-integral}
\end{align}
So:
\begin{align}
c = \frac{m_{\mathrm{L}}}{Q} \label{eq:model-proportionality-constant}
\end{align}

\section{Estimating the uncertainty in \textit{c}}
The total uncertainty is given by:
\begin{align}
\delta c = \left[ \left( \frac{\delta m_{\mathrm{L}}}{m_{\mathrm{L}}} \right)^2 + \left(\frac{\delta Q}{Q}\right)^2 \right]^{1/2} \label{eq:uncertainty-c_implicit}
\end{align}

\subsection{\textit{Q} is estimated numerically}
$Q$ is estimated using either the trapezoid or Simpson's rule.

\textbf{Simpson's rule}
\begin{align}
Q = \frac{\Delta t}{3} \sum_{i=0}^N w_i, I_i T_{e,i}, \qquad w_i \in \{1,4,2,4,2,\ldots,4,1\}
\label{eq:Q-simpson}
\end{align}
So each contribution is $q_i = \frac{\Delta t}{3}w_i I_i, T_{e,i}$:
\begin{align*}
Q = \sum_{i=0}^N q_i
\end{align*}

The partial derivatives of $q_i$ are
\begin{align*}
\frac{\partial q_i}{\partial I_i} = \frac{q_i}{I_i}, \quad \frac{\partial q_i}{\partial T_{e,i}} = \frac{q_i}{T_{e,i}}
\end{align*}
So
\begin{align*}
\delta q_i &= \left[ \left( \frac{q_i \delta I_i^2}{I_i} \right)^2 + \left( \frac{q_i \delta T_{e,i}^2}{T_{e,i}} \right)^2 \right]^{1/2} \\
&= q_i\left[ \left( \frac{\delta I_i^2}{I_i} \right)^2 + \left( \frac{\delta T_{e,i}^2}{T_{e,i}} \right)^2 \right]^{1/2}
\end{align*}

The uncertainty in $Q$ is
\begin{align*}
\delta Q = \left[\sum_i \left(\frac{\partial Q}{\partial q_i}\right)^2 \delta q_i^2\right]^{1/2} = \left[\sum_i  \delta q_i^2\right]^{1/2}
\end{align*}

Then,
\begin{align*}
\delta Q = \left\{\sum_i q_i^2 \left[  \left( \frac{\delta I_i}{I_i}\right)^2 + \left( \frac{\delta T_{e,i}}{T_{e,i}}\right)^2\right]\right\}^{1/2}
\end{align*}

Expanding $q_i$:
\begin{align*}
\delta Q &=
\left\{
  \sum_i \left( \frac{\Delta t}{3} w_i I_i T_{e,i} \right)^2  \left( \frac{\delta I_i}{I_i}\right)^2
  + \sum_i \left( \frac{\Delta t}{3} w_i I_i T_{e,i} \right)^2 \left( \frac{\delta T_{e,i}}{T_{e,i}}\right)^2
\right\}^{1/2}
\end{align*}

Considering the systematic uncertainty on the current: $\delta I_i / I_i \equiv \epsilon_I$ (from considering a $\epsilon_I = $ 5\% tolerance resistor and $I = V/R \Rightarrow \delta I = I \delta R / R = \epsilon_I I$ ), and canceling terms:
\begin{align*}
\delta Q &=
\left\{
  \sum_i \left( \frac{\Delta t}{3} w_i T_{e,i} I_i \epsilon_I \right)^2  
  + \sum_i \left( \frac{\Delta t}{3} w_i I_i  \delta T_{e,i} \right)^2
\right\}^{1/2}
\end{align*}

The first term is equal to $\epsilon_i Q$:
\begin{align}
\boxed{
\delta Q =
\left\{
  \left(\epsilon_I Q\right)^2
  + \sum_i \left( \frac{\Delta t}{3} w_i I_i  \delta T_{e,i} \right)^2 \label{eq:uncertainty-Q}
\right\}^{1/2}
}
\end{align}

This can be implemented in Python as follows:
\begin{lstlisting}[caption={Example implementation of $\delta Q$ in Python}]
import numpy as np
from scipy.integrate import simpson

# --- Inputs ---
# t, I, Te, dTe are 1D arrays of equal length N 
# (N-1 must be even for pure Simpson)
# ddm = uncertainty in delta_m, dm = delta_m value
epsilon_I = 0.05  # delta_R / R

# --- Integral via Simpson's rule ---
integrand = I * Te
Q = simpson(integrand, x=t)

# --- Build Simpson weights explicitly ---
N = len(t)
dt = t[1] - t[0]  # assumes uniform spacing

def simpson_weights(N, dt):
    """Returns Simpson's 1/3 rule weights for N points (N must be odd)."""
    assert N % 2 == 1, "Need odd number of points (even number of intervals)"
    w = np.ones(N)
    w[1:-1:2] = 4   # odd indices
    w[2:-2:2] = 2   # even indices (interior)
    return w * dt / 3

w = simpson_weights(N, dt)

# --- Uncertainty in Q ---
dQ_sys  = epsilon_I * Q                      # systematic (R calibration)
dQ_stat = np.sqrt(np.sum((w * I * dTe)**2))  # statistical (Te errors)
dQ      = np.sqrt(dQ_sys**2 + dQ_stat**2)

# --- Uncertainty in c ---
c = dm / Q
rel_unc_c = np.sqrt((ddm / dm)**2 + (dQ / Q)**2)
delta_c   = c * rel_unc_c

print(f"Q        = {Q:.4e}")
print(f"dQ (sys) = {dQ_sys:.4e}  ({100*dQ_sys/Q:.1f}%)")
print(f"dQ (stat)= {dQ_stat:.4e}  ({100*dQ_stat/Q:.1f}%)")
print(f"dQ total = {dQ:.4e}  ({100*dQ/Q:.1f}%)")
print(f"c        = {c:.4e}")
print(f"delta_c  = {delta_c:.4e}  ({100*rel_unc_c:.1f}%)")
\end{lstlisting}

\section{Expanding the uncertainty in the rate of mass loss}
The rate of mass loss is given by Eq.~\eqref{eq:mass-loss-I}:
\begin{align}
r(t) = \frac{dm}{dt} = c I(t) T_e(t). \label{eq:mass-loss-rate}
\end{align}
The constant $c$, the current $I(t)$, and the electron temperature $T_e(t)$ have associated uncertainties $\delta c$, $\delta I = \epsilon_I $, and $\delta T_e$.

Writing Eq.~\eqref{eq:mass-loss-rate} explicitly
\begin{align*}
r(t) = \frac{m_{\mathrm{L}}}{\int_0^{\tau}I(t)T_e(t) dt} I(t) T_e(t). 
\end{align*}
Take the logarithm
\begin{align}
\ln r(t) &= \ln \left(m_{\mathrm{L}}\right) + \ln \left[ I(t) \right] + \ln \left[T_e(t)\right] - \ln \left( \int_0^{\tau}I(t)T_e(t) dt\right) \notag\\
&= \ln m_{\mathrm{L}} + \ln  I + \ln T_e - \ln Q
\end{align}

Take the partial derivative
\begin{align}
\frac{\partial \ln r}{\partial m_{\mathrm{L}}} &= \frac{1}{r} \frac{\partial r}{\partial m_{\mathrm{L}}} = \frac{1}{m_{\mathrm{L}}}
\end{align}

Varying each error source, \textbf{one at a time}:

\subsection{ Variation 1: $\delta m_{\mathrm{L}}$ (independent)}
\begin{align}
\left(\frac{\delta r}{r}\right)_{m_{\mathrm{L}}} = \frac{1}{r}\left(\frac{\partial r}{\partial m_{\mathrm{L}}}\right)\delta m_{\mathrm{L}} = \frac{\delta m_{\mathrm{L}}}{m_{\mathrm{L}}}.\tag{8}
\end{align}

\subsection{Variation 2: $\delta R/R$ (systematic, affects all $I_i$ and $I(t)$}
Every $I_i$ has propagates the uncertainty as $I_i \to I_i(1+\epsilon)$ and $I(t)\to I(t)(1+\epsilon)$. 

We are assuming that the error is systematic: \textbf{$R$ has a fractional error $\epsilon = \delta R /R$}.

Since $I = V/R$, a resistor that is \textbf{larger than assumed} means the \textbf{current is smaller than assumed}:
\begin{align*}
I_{\mathrm{true}} = \frac{V}{R_{\mathrm{true}}} = \frac{V}{R(1+\epsilon)}\approx I(1-\epsilon)
\end{align*}

So strictly speaking $\delta I / I = -\epsilon$. However, for \textbf{uncertainty propagation, we only care about the magnitude} of the effect, not the sign. Writing $I \to I (1+\epsilon)$, with $\epsilon = \pm \delta R/R$ is just a convenient way to track how a fractional shift propagates through the formula.

Then, $Q \to Q(1+\epsilon)$:
\begin{align*}
Q = \int_0^{\tau} I(t) T_e(t) dt \to \int_0^{\tau} I(t) (1 + \epsilon) T_e(t) dt
= (1 + \epsilon)\int_0^{\tau} I(t)  T_e(t) dt = Q (1 + \epsilon).
\end{align*}

\textbf{This causes cancelation in $r$}:
Substitute both shifts in $r$:
\begin{align*}
r &= c I(t) T_e(t) = \frac{m_{\mathrm{L}}}{Q} I(t)T_e(t) \\
r \to \frac{m_{\mathrm{L}}}{Q(1+\epsilon)} I(t)(1+\epsilon)T_e(t)
&= \frac{m_{\mathrm{L}}}{Q} I(t)T_e(t) \frac{(1+\epsilon)}{(1+\epsilon)} = r
\end{align*}

The $(1+\epsilon)$ factors cancel \textbf{exactly}, so $r$ is completely insensitive to $\epsilon$

\subsection{Variation 3: $\delta T_{e}$}
Assuming Simpson's rule, $Q = \frac{\Delta t}{3}\sum_{i}w_iI_iT_{e,i}$. Then,
\begin{align}
\ln r_k = \ln m_{\mathrm{L}} + \ln I_k + ln T_{e,k} - \ln \left(\frac{\Delta t}{3}\sum_{j}w_jI_jT_{e,j} \right)
\end{align}

$T_e$ appears \textbf{twice}:
\begin{itemize}
    \item Once \textbf{explicitly} in $+\ln T_{e,l}$
    \item Once \textbf{implicitly} inside $-\ln Q$, at every time step $j$, including $j=k$
\end{itemize}

\textbf{Case $j \neq k$: $T_{e,j}$ appears only inside $\ln Q$}
\begin{align}
\frac{\partial \ln r_k}{\partial T_{e,j}} = -\frac{\partial \ln Q}{\partial T_{e,j}} = -\frac{1}{Q}\frac{\partial Q}{\partial T_{e,j}} = - \frac{w_j I_j \frac{\Delta t}{3}}{Q} \tag{9}
\end{align}

\textbf{Case $j = k$: $T_{e,j}$ appears in both terms $\ln T_{e,k}$, and $-\ln Q$ simultaneously}
\begin{align}
\frac{\partial \ln r_k}{\partial T_{e,k}} &= \underbrace{\frac{\partial \ln T_{e,k}}{\partial T_{e,k}}}_{\text{explicit term}} -\underbrace{\frac{\partial \ln Q}{\partial T_{e,k}}}_{\text{implicit term}} \notag \\
&= \frac{1}{T_{e,k}} - \frac{w_j I_j \frac{\Delta t}{3}}{Q}
\end{align}

Multiplying through by $\delta T_{e,k}$ to express as a relative uncertainty:
\begin{align}
\frac{\partial \ln r_k}{\partial T_{e,k}} \delta T_{e,k} = \left[ 1 - \frac{w_k I_k T_{e,k} \frac{\Delta t}{3}}{Q}\right]\frac{\delta T_{e,k}}{T_{e,k}}
\end{align}

Define
\begin{align}
\boxed{
\alpha_k \equiv \frac{w_k I_k T_{e,k} \frac{\Delta t}{3}}{Q} \label{eq:alpha_k}
}
\end{align}

\subsection{Full $T_e$ uncertainty}
We can unify both cases cleanly
\begin{align}
\left(\frac{\delta r_k}{r_k}\right)_{T_e}^2 &= \sum_{\color{red}{j}}\left( \frac{\partial \ln r_k}{\partial T_{e,j}} \delta T_{e,j}\right)^2
=
\sum_{\color{red}{j}}\left( \frac{\partial }{\partial T_{e,j}} \left[ \ln T_{e,k} - \ln Q \right]\delta T_{e,j}\right)^2 
\end{align}
where
\begin{align}
\frac{\partial \ln r_k}{\partial T_{e,j}} =
\begin{cases}
\dfrac{1-\alpha_k}{T_{e,k}}, & j=k\\[1em]
-\dfrac{w_j I_j \frac{\Delta t}{3}}{Q}, & j\neq k
\end{cases}
\end{align}

So the sum over \textbf{all} $j$ naturally splits into the two cases:
\begin{align}
\boxed{
\left(\frac{\delta r_k}{r_k}\right)^2 = (1 - \alpha_k)^2\left(\frac{\delta T_{e,k}}{T_{e,k}}\right)^2 + \sum_{j\neq k}\left(\frac{w_jI_j\frac{\Delta t}{3}}{Q}\right)^2  \delta T_{e,j}
}
\end{align}

This can be easily implemented in \textbf{Python}
\begin{lstlisting}[caption={Sample implementation of $\delta r_k$ in Python}]
import numpy as np

# --- Trapezoid weights ---
def trapezoid_weights(t):
    """Non-uniform spacing supported."""
    dt = np.diff(t)
    w = np.zeros(len(t))
    w[:-1] += dt / 2
    w[1:]  += dt / 2
    return w

w = trapezoid_weights(t)

# --- Integral ---
Q = np.sum(w * I * Te)

# --- Alpha: fractional contribution of each time step ---
alpha = w * I * Te / Q  # shape (N,), sums to 1

# --- Uncertainty in r(t) ---
r = c * I * Te

# Term 1: delta_m (scalar)
frac_dm = (ddm / dm)**2

# Terms 2+3: T_e contributions, split at each t
# For each index k, the local term is (1 - alpha[k])^2 * (dTe[k]/Te[k])^2
# and the sum over j != k of (w[j]*I[j]*dTe[j]/Q)^2

full_Te_sum = np.sum((w * I * dTe / Q)**2)  # sum over ALL j
local_Te    = (w * I * dTe / Q)**2          # contribution of each j 
                                            # to sum

# For index k:
#   sum_{j != k} = full_Te_sum - local_Te[k]
#   local term   = (1 - alpha[k])^2 * (dTe[k]/Te[k])^2

frac_Te_nonlocal = full_Te_sum - local_Te         # array, sum over j!=t
frac_Te_local    = (1 - alpha)**2 * (dTe / Te)**2 # array, local t term

delta_r_over_r = np.sqrt(frac_dm + frac_Te_local + frac_Te_nonlocal)
delta_r = r * delta_r_over_r
\end{lstlisting}

\section{Propagating uncertainty in $m_{\mathrm{lost}}(t)$}
The cumulative eroded mass up to time step $k$ is:
\begin{align}
    m_k = c S_k = m_{\mathrm{L}}\left(\frac{S_k}{S}\right),
    \label{eq:mass_loss_equation}
\end{align}
where $S_k \equiv \sum_{i=0}^k w_i, I_i, T_{e,i}$, and $S \equiv \sum_{i=0}^N w_i, I_i, T_{e,i} = Q / (\Delta t / 3)$. For simplicity, we designate $m$ as the emitted (or lost) mass $m_{\mathrm{lost}}$.

Taking the log:
\begin{align}
    \ln m_k = \ln m_{\mathrm{L}} + \ln S_k - \ln S
\end{align}
\subsection{Variations}
\subsubsection{$\delta m_{\mathrm{L}}$ (independent)}
\begin{align}
    \frac{\partial \ln m_k}{\partial m_{\mathrm{L}}} = \frac{1}{m_{\mathrm{L}}}
    \Longrightarrow 
    \left( \frac{\delta m_k}{m_k}\right)_{m_{\mathrm{L}}} = \frac{\delta m_{\mathrm{L}}}{m_{\mathrm{L}}}
\end{align}

\subsubsection{$\delta R/R$ (systematic, all $I_i$ scale together}
Every $I_i \to I_i(1+\epsilon)$, so $S_k\to S_k (1+ \epsilon)$, and $S \to S (1+\epsilon)$.
\begin{align*}
    \delta \ln m_k &= \ln \left[S_k (1 + \varepsilon) \right] -  \ln \left[S (1 + \varepsilon) \right]
    - \ln S_k + \ln S \\
    &= \ln (1+\epsilon) - \ln (1+\epsilon) = 0
\end{align*}
\textbf{Cancels exactly}, same argument as for $r_k$.

\subsubsection{$\delta T_{e,j}$ -- Three cases now}
Since $\ln m_k = \ln m_{\mathrm{L}} + \ln S_k - \ln S$, varying $T_{e,j}$, gives
\begin{align}
    \frac{\partial \ln m_k}{\partial T_{e,j}} = \frac{\partial \ln S_k}{\partial T_{e,j}} - \frac{\partial \ln S}{\partial T_{e,j}}
\end{align}
The key difference from $r_k$: there is \textbf{no explicit} $\ln T_{e, k}$ \textbf{term}. $T_e$ only appears inside the two $S_k$ and $S$. \emph{Why do we need to vary over all values of $j$?}: Because $m_k$ depends on \textbf{all} $T_{e,j}$ values, not just the local one at $j=k$: Look at $m_k = m_{\mathrm{L}} (S_k/S) = m_{\mathrm{L}}\left(\sum_{i=0}^k w_i I_i T_{e,j}\right)/\left(\sum_{i=0}^N w_i I_i T_{e,j}\right)$. Every single $T_{e,j}$ for $j=0,\ldots,N$ appears somewhere in this expression.

\begin{description}
    \item[Case 1: $j > k$ -- appears only in $S$, not in $S_k$]
    \begin{align*}
        \frac{\partial \ln m_k}{\partial T_{e,j}} &= 
        \frac{1}{S_k}\frac{\partial S_k}{\partial T_{e,j}} - \frac{1}{S}\frac{\partial S}{\partial T_{e,j}} \\
        &=
        \frac{1}{S_k}\sum_{i=0}^k \underbrace{\frac{\partial }{\partial T_{e,j}} \left(w_i I_i T_{e,i}\right)}_{=0 \text{ since } j > k} - \frac{1}{S} \sum_{i=0}^N\underbrace{\frac{\partial }{\partial T_{e,j}} \left(w_i I_i T_{e,i}\right)}_{=\delta_{ij}} \\
        &= -\frac{w_j I_j }{S} \\
        \Rightarrow \frac{\partial \ln m_k}{\partial T_{e,j}}\delta T_{e,j} &= -\frac{w_j I_j}{S} \delta T_{e,j}
    \end{align*}
    \item[Case 2: $j < k$ -- appears in both $S_k$, and $S$]
    \begin{align*}
        \frac{\partial \ln m_k}{\partial T_{e,j}} &= \frac{1}{S_k}\sum_{i=0}^k \underbrace{\frac{\partial }{\partial T_{e,j}} \left(w_i I_i T_{e,i}\right)}_{=\delta_{ij}} - \frac{1}{S}\sum_{i=0}^N\underbrace{\frac{\partial }{\partial T_{e,j}} \left(w_i I_i T_{e,i}\right)}_{=\delta_{ij}} \\
        &=
        \frac{w_j I_j}{S_k} - \frac{w_j I_j }{S} = w_j I_j  \left( \frac{1}{S_k}  - \frac{1}{S}\right) \\
        \Rightarrow \frac{\partial \ln m_k}{\partial T_{e,j}} \delta T_{e,j} &=
        w_j I_j \left( \frac{1}{S_k}  - \frac{1}{S}\right) \delta T_{e,j}
    \end{align*}
    \item[Case 3: $j=k$ -- also  appears in both $S_k$, and $S$] 
    \begin{align*}
        \frac{\partial \ln m_k}{\partial T_{e,k}} &= \frac{1}{S_k}\sum_{i=0}^k \underbrace{\frac{\partial }{\partial T_{e,k}} \left(w_i I_i T_{e,i}\right)}_{=\delta_{ik}} - \frac{1}{S}\sum_{i=0}^N\underbrace{\frac{\partial }{\partial T_{e,k}} \left(w_i I_i T_{e,i}\right)}_{=\delta_{ik}} \\
        &=
        \frac{w_k I_k}{S_k} - \frac{w_k I_k }{S}
        =
        w_k I_k \left( \frac{1}{S_k}  - \frac{1}{S}\right) \\
        \Rightarrow \frac{\partial \ln m_k}{\partial T_{e,k}} \delta T_{e,k} 
        &=
        w_k I_k \left( \frac{1}{S_k}  - \frac{1}{S}\right) \delta T_{e,k} 
    \end{align*}
\end{description}
The full expression for $\frac{\partial \ln m_k}{\partial T_{e,j}}\delta T_{e,j}$ is
\begin{align}
    \frac{\partial \ln m_k}{\partial T_{e,k}} \delta T_{e,j} =
    \begin{cases}
        w_j I_j \left( \dfrac{1}{S_k} - \dfrac{1}{S}\right)\delta T_{e,j} & j \leq k \\[1.5em]
        -\dfrac{w_jI_j}{S}\delta T_{e,j} & j > k
    \end{cases}
\end{align}
\subsubsection{Full result}
\begin{align}
    \boxed{
    \frac{\delta m_k}{m_k} = \left\{
        \left(\frac{\delta m_{\mathrm{L}}}{m_{\mathrm{L}}}\right)^2
        +
        \sum_{j\leq k} \left[w_j I_j \left( \frac{1}{S_k} - \frac{1}{S}\right)\delta T_{e,j}\right]^2
        +
        \sum_{j > k} \left[ \frac{w_j I_j}{S}\delta T_{e,j}\right]^2
    \right\}^{1/2} 
    }
\end{align}

This can be easily implemented in \textbf{Python}:
\begin{lstlisting}[caption={Example code to estimate $\delta m_k$ in Python}]
# S_k: cumulative sum up to each index k
S_k = np.cumsum(w * I * Te)   # shape (N,)
S   = S_k[-1]                 # full sum = Q

# Precompute per-step weight * current
wI = w * I                    # shape (N,)

# For each k, build delta_m_k/m_k
frac_dm = (ddm / dm)**2       # scalar

frac_Te = np.zeros(N)
for k in range(N):
    past   = np.arange(0, k+1)          # j <= k
    future = np.arange(k+1, N)          # j > k
    term_past   = np.sum((wI[past]   * (1/S_k[k] - 1/S) * dTe[past])**2)
    term_future = np.sum((wI[future] * dTe[future] / S)**2)
    frac_Te[k]  = term_past + term_future

delta_m_over_m = np.sqrt(frac_dm + frac_Te)
\end{lstlisting}

\section{Change of $A_w$ with incident heat load}
We assume that the wetted area $A_w$ decreases with exposure to heat load. This should be proportional to the mass loss:
\begin{align*}
    A_{\mathrm{w}} \equiv Dh(t)
\end{align*}
where $D$ is the diameter of the rod and $h(t)$ is the height of the rod exposed to the plasma at time $t_k$. $h(t)$ decreases with mass loss: 
\begin{align*}
    \Delta h(t) \propto m_{\mathrm{lost}}(t) 
\end{align*}
with $m_{\mathrm{lost}}$ as defined in Eq.~\eqref{eq:mass_loss_equation}. We assume $D$ remains constant.

The volume of the rod lost during exposure is $V(t) = \pi D^2 h(t) / 4$. Assuming a constant effective density $\rho$, $m_{\mathrm{lost}}(t)/\rho = \pi D^2 \Delta h(t) / 4$. Then,
\begin{align*}
    \Delta h(t) &= \frac{4}{\rho\pi D^2} m_{\mathrm{lost}}(t)
\end{align*}

And the wetted area is given by
\begin{align}
    A_w(t) = D\left( \frac{4}{\rho_{\mathrm{rod}}\pi D^2}\right) m_{\mathrm{lost}}(t) 
    = \frac{4}{\rho_{\mathrm{rod}}\pi D} m_{\mathrm{lost}} (t) 
    \label{eq:A_w_t}
\end{align}

\subsection{Uncertainty propagation in $A_w$}
The quantities $\rho_{\mathrm{rod}}$, $D$, and $m_{\mathrm{lost}}$ have associated uncertainties $\delta \rho_{\mathrm{rod}}$, $\delta D$, and $\delta m_{\mathrm{lost}}$. So the uncertainty in $A_w$ is given by
\begin{align}
    \delta A_w (t) = A_w(t) \left\{ 
        \left(\frac{\delta D}{D}\right)^2 
        +
        \left(\frac{\delta \rho_{\mathrm{rod}}}{\rho_{\mathrm{rod}}}\right)^2 
        + 
        \left(\frac{\delta m_{\mathrm{lost}}}{m_{\mathrm{lost}}}\right)^2
    \right\}^{1/2}
\end{align}

\subsection{Surface recession rate}
The surface recession rate for the wetted surface 
\begin{align}
    \nu \equiv \frac{1}{\rho_{\mathrm{rod}} A_w(t)} \frac{dm(t)}{dt} 
\end{align}
Plugging in Eq.~\eqref{eq:A_w_t}:
\begin{align*}
    \nu(t) = \frac{1}{\rho_{\mathrm{rod}}} \left( \frac{\rho_{\mathrm{rod} \pi D}}{4 m_{\mathrm{lost}}(t)}\right) r(t)
    = \frac{\pi D}{4} \frac{r(t)}{m_{\mathrm{lost}}(t)}
\end{align*}
Which can be written \textbf{explicitly} in terms of $I(t)$, and $T_e(t)$ as
\begin{align}
    \nu(t) = \frac{\pi D}{4} \frac{1}{c \int_0^t I(t')T_e(t') dt} c I(t) T_e(t) 
    \label{eq:nu_explicitly}
\end{align}


Expressing it in numeric form:
\begin{align}
    \nu_k = \frac{\pi D}{4} \frac{I_k T_{e,k}}{Q_k} 
    \label{eq:surface_recession_Aw}
\end{align}

\subsection{Uncertainty in surface recession rate}
Start with the expression in Eq.~\eqref{eq:surface_recession_Aw}:
\begin{align*}
    \nu_k = \frac{\pi D}{4} \frac{I_k T_{e,k}}{Q_k}.
\end{align*}


Define
\begin{align}
    X_i \equiv I_i T_{e,i}, \quad Q_k \equiv \sum_{i=0}^k w_i X_i, 
\end{align}
where we have lumped $\Delta t/3$ in the quadrature weights $w_i$.

The propagation of the uncertainty of $f(u_1,\ldots,u_n)$ is given by
\begin{align*}
    \sigma_f^2 = \sum_m \left(\frac{\partial f}{\partial u_m}\right)^2\sigma^2_m
\end{align*}

Assume $\delta V = 0$, so that the error in $I=V/R$ is systematic from $\delta R/R =\epsilon$. Because $R$ is constant across measurements, \textbf{even with a nonzero $\epsilon$, it does not contribute to $\delta_k$ at all.}

Here is why, if $Ij = V_j/R_j$ with the \emph{same} $R$ is used for every sample, then substitute it directly into $\nu_k$ before ever differentiating:
\begin{align*}
    \nu_k = \frac{\pi D}{4}\frac{I_k T_{e,k}}{Q_k} = \frac{\pi D}{4}\frac{(V_k/R)T_{e,k}}{\sum_{i}^kw_i(V_i/R)T_{e,i}} =  \frac{\pi D}{4} \frac{V_k T_e{,k}}{\sum_{i}^kw_iV_i T_{e,i}}
\end{align*}
$R$ factors out of every term in the sum $S_k$ and out of the numerator identically, so it cancels completely --algebraically, before taking any derivative.

\subsubsection{The $D$ partial}
$\nu_k \propto D$ exactly, so
\begin{align}
    \frac{\partial \nu_k}{\partial D} = \frac{\nu_k}{D}
    \label{eq:partial_vk_D}
\end{align}
trivial.

\subsubsection{The $X_j$ partial}
Rather than differentiating directly with respect to $T_{e,j}$, it is cleaner to first get $\partial \nu_k/\partial X_j$, then use the chain rule through $X_j = I_j T_{e,j}$.

Write $\nu_k = \frac{\pi D}{2} X_k Q_k^{-1}$ and use the product rule
\begin{align*}
    \frac{\partial \nu_k}{\partial X_j} &= \frac{\pi D}{4}\left[ \frac{\partial X_k}{\partial X_j}\cdot \frac{1}{Q_k} + X_k\cdot\frac{\partial \left(Q_k\right)^{-1}}{\partial X_j}\right]
\end{align*}

Two pieces to work out:
\begin{itemize}
    \item $\dfrac{\partial X_k}{\partial X_j} = \delta_{jk}$ (the Kronecker delta -- $X_k$ only depends on itself, so this is 1 if $j=k$, 0 otherwise). \textbf{This is the term that captures the numerator/denominator correlation} -- $X_k$ appears explicitly in the numerator , so differentiating the numerator with respect to $X_j$ only ``sees'' it when $j=k$.
    \item $\dfrac{\partial Q_k}{\partial X_j} = w_j$ for $j\leq k$ (since $Q_k=\sum_{i=0}^k w_i X_i$ is linear in each $X_i$), so by the chain rule, $\dfrac{\partial \left(S_k\right)^{-1}}{\partial X_j} = - Q_k^{-2}\dfrac{\partial Q_k}{\partial X_j} = -\dfrac{w_j}{Q_k^2}$
\end{itemize}

Putting these together:
\begin{align}
    \frac{\partial \nu_k}{\partial X_j} &= \frac{\pi D}{4}\left[ \frac{\delta_{jk}}{Q_k} - \frac{w_j X_k}{Q_k^2} \right], \quad j\leq k
    \label{eq:partial-vk_Xj}
\end{align}

\subsubsection{Chain rule down to $T_{e,j}$}
Since $X_j = I_jT_{e,j}$, with no error through $I_j$ (systematic error $\delta R/R$ cancels out and $\delta V_j = 0$ assumed),
\begin{align*}
    \frac{\partial v_k}{\partial T_{e,j}} = \frac{\partial v_k}{\partial X_j}\cdot \frac{\partial X_j}{\partial T_{e,j}} = \frac{\partial v_k}{\partial X_j} \cdot I_j
\end{align*}
so in variance form
\begin{align*}
    \frac{\partial v_k}{\partial T_{e,j}} \delta T_{e,j} = \frac{\pi D}{4}\left[ \frac{\delta_{jk}}{Q_k} - \frac{w_j X_k}{Q_k^2} \right] I_j \delta T_{e,j}, \quad j\leq k
\end{align*}

It is cleanest to just define $\delta X_j \equiv I_j \delta T_{e,j}$ and treat $\left(\frac{\partial v_k}{\partial X_j} \right)^2 \delta X_j^2$ as the contribution.

\subsubsection{Assemble the total variance}
Sum the independent contributions from $D$ and each $T_{e,j}, j=0,1,\ldots,k$:
\begin{align}
    \boxed{
        \delta v_k = \left\{\left( \frac{v_k}{D}\right)^2 \delta D^2 + \sum_{j=0}^k\left[ \frac{\pi D}{4} \left( \frac{\delta_{jk}}{Q_k} - \frac{w_j X_k}{Q_k^2}\right)\right]^2\delta X_j^2\right\}^{1/2}
    }
\end{align}

This can be implemented in Python as
\begin{lstlisting}[caption={Example implementation of the recession rate calculation in Python}]
"""
Uncertainty propagation for the DiMES surface recession rate.

    nu_k = (pi * D / 2) * X_k / Q_k,   X_i = I_i * Te_i
    Q_k  = numerical integral of X over [t_0, t_k], via 'trapezoid' or
           'simpson' composite quadrature -- NON-UNIFORM time grid t.
    I_i treated as exact (dV = 0); dR/R contributes nothing (see prior
    derivation) and is not modeled.

Quadrature coefficients dQ_k/dX_j  (t is now an array, not a scalar dt)
-------------------------------------------------------------------------
trapezoid:  dQ_k/dX_0 = (t_1-t_0)/2
            dQ_k/dX_j = (t_{j+1}-t_{j-1})/2   (0<j<k)
            dQ_k/dX_k = (t_k-t_{k-1})/2

simpson:    composite Simpson over PANELS of 3 consecutive points, using
            the unequal-interval formula (reduces to the standard
            h/3, 4h/3, h/3 weights when the panel's two sub-intervals are
            equal). For odd k, a plain trapezoidal correction (using the
            true local width t_k - t_{k-1}) is added for the leftover
            interval, exactly as in the uniform-grid version.
"""

import numpy as np


def quad_coeffs(k, t, method="trapezoid"):
    """dQ_k/dX_j for j = 0..k, for the chosen quadrature method.
    t : array_like of sample times (length >= k+1), need not be uniform."""
    t = np.asarray(t, dtype=float)
    c = np.zeros(k + 1)
    if k == 0:
        return c

    if method == "trapezoid":
        c[0] = (t[1] - t[0]) / 2.0
        c[1:k] = (t[2:k + 1] - t[0:k - 1]) / 2.0
        c[k] = (t[k] - t[k - 1]) / 2.0
        return c

    if method == "simpson":
        m = k if k % 2 == 0 else k - 1  # largest even index <= k
        for start in range(0, m, 2):
            h0 = t[start + 1] - t[start]
            h1 = t[start + 2] - t[start + 1]
            w0 = (h0 + h1) / 6.0 * (2.0 - h1 / h0)
            w1 = (h0 + h1) / 6.0 * (h0 + h1) ** 2 / (h0 * h1)
            w2 = (h0 + h1) / 6.0 * (2.0 - h0 / h1)
            c[start] += w0
            c[start + 1] += w1
            c[start + 2] += w2
        if k % 2 == 1:  # trailing trapezoid correction, true local width
            half = (t[k] - t[k - 1]) / 2.0
            c[k - 1] += half
            c[k] += half
        return c

    raise ValueError(f"unknown method: {method}")


def cumulative_integral(X, t, method="trapezoid"):
    """Q_k for k = 0..N-1 (Q_0 = 0)."""
    N = len(X)
    Q = np.zeros(N)
    for k in range(1, N):
        c = quad_coeffs(k, t, method)
        Q[k] = np.dot(c, X[: k + 1])
    return Q


def propagate_analytic(D, dD, I, Te, dTe, t, method="trapezoid"):
    """
    Parameters
    ----------
    D, dD : float
    I : array_like, shape (N,)      -- treated as exact (dV = 0)
    Te, dTe : array_like, shape (N,)
    t : array_like, shape (N,)      -- sample times, non-uniform OK
    method : 'trapezoid' or 'simpson'

    Returns
    -------
    nu, dnu : ndarray, shape (N,)   (index 0 is undefined -- no interval yet)
    """
    I, Te, dTe, t = (np.asarray(a, dtype=float) for a in (I, Te, dTe, t))
    N = len(I)

    X = I * Te
    dX = I * dTe
    Q = cumulative_integral(X, t, method)

    nu = np.full(N, np.nan)
    dnu = np.full(N, np.nan)
    for k in range(1, N):
        Qk = Q[k]
        nu[k] = (np.pi * D / 4.0) * X[k] / Qk

        c = quad_coeffs(k, t, method)          # dQ_k/dX_j, j=0..k
        j = np.arange(k + 1)
        kronecker = (j == k).astype(float)
        dnu_dXj = (np.pi * D / 4.0) * (kronecker / Qk - c * X[k] / Qk ** 2)
        var_X = np.sum((dnu_dXj * dX[: k + 1]) ** 2)
        var_D = (nu[k] / D * dD) ** 2
        dnu[k] = np.sqrt(var_X + var_D)

    return nu, dnu


def propagate_montecarlo(D, dD, I, Te, dTe, t, method="trapezoid",
                          n_draws=20000, seed=0):
    rng = np.random.default_rng(seed)
    I, Te, dTe, t = (np.asarray(a, dtype=float) for a in (I, Te, dTe, t))
    N = len(I)

    D_s = rng.normal(D, dD, size=n_draws)
    Te_s = rng.normal(Te, dTe, size=(n_draws, N))
    X_s = I[None, :] * Te_s

    Q_s = np.zeros((n_draws, N))
    for k in range(1, N):
        c = quad_coeffs(k, t, method)
        Q_s[:, k] = X_s[:, : k + 1] @ c

    nu_s = np.full((n_draws, N), np.nan)
    nu_s[:, 1:] = (np.pi * D_s[:, None] / 2.0) * X_s[:, 1:] / Q_s[:, 1:]

    return np.nanmean(nu_s, axis=0), np.nanstd(nu_s, axis=0)


if __name__ == "__main__":
    N = 51
    rng = np.random.default_rng(1)
    # non-uniform grid: irregular but increasing spacings
    steps = 0.5 + rng.random(N - 1)
    t = np.concatenate([[0.0], np.cumsum(steps)])
    t /= t[-1]  # normalize to [0,1]

    I = 5.0 + 3.0 * np.sin(np.pi * t)
    Te = 20.0 + 10.0 * t
    dTe = 0.10 * Te
    D, dD = 0.6, 0.02

    for method in ("trapezoid", "simpson"):
        print(f"\n=== {method} (non-uniform grid) ===")

        X = I * Te
        k_test = 21  # odd, exercises the boundary-correction branch
        eps = 1e-6
        Q0 = cumulative_integral(X, t, method)
        analytic_c = quad_coeffs(k_test, t, method)
        fd_c = np.zeros(k_test + 1)
        for j in range(k_test + 1):
            Xp = X.copy(); Xp[j] += eps
            Qp = cumulative_integral(Xp, t, method)
            fd_c[j] = (Qp[k_test] - Q0[k_test]) / eps
        print("max |analytic dQ_k/dX_j - finite-diff| =",
              np.max(np.abs(analytic_c - fd_c)))

        nu, dnu = propagate_analytic(D, dD, I, Te, dTe, t, method)
        nu_mc, dnu_mc = propagate_montecarlo(D, dD, I, Te, dTe, t, method)

        print(f"{'k':>3} {'nu':>10} {'dnu (analytic)':>16} {'dnu (MC)':>10}")
        for k in (1, N // 4, N // 2, 3 * N // 4, N - 1):
            print(f"{k:3d} {nu[k]:10.4f} {dnu[k]:16.4f} {dnu_mc[k]:10.4f}")
\end{lstlisting}
    
\end{document}
